% ==============================================================================
% DOCUMENTO TÉCNICO V2 - REQUERIMIENTOS Y DISEÑO ARQUITECTÓNICO
% Proyecto: Sistema Integral de Gestión de Capacitaciones (SIGC-CUSCO)
% Asignatura: Ingeniería de Software I - UNSAAC
% ==============================================================================

\documentclass[11pt,a4paper]{article}

\usepackage[utf8]{inputenc}
\usepackage[T1]{fontenc}
\usepackage[spanish, es-tabla, es-nodecimaldot]{babel}
\spanishdeactivate{"}
\usepackage{csquotes}

% Márgenes formales
\usepackage[margin=2.2cm, headheight=16pt]{geometry}
\usepackage{mathptmx}
\usepackage{setspace}
\onehalfspacing

\usepackage{amsmath,amssymb,amsfonts}
\usepackage{graphicx}
\usepackage{float}
\usepackage{booktabs}
\usepackage{tabularx}
\usepackage{multirow}
\usepackage{array}
\usepackage{pdflscape}
\usepackage{caption}
\usepackage{subcaption}

% Carga previa de xcolor con opciones antes de tcolorbox para evitar warnings
\usepackage[table,xcdraw]{xcolor}
\usepackage[most]{tcolorbox}

\usepackage{tikz}
\usetikzlibrary{arrows.meta, positioning, shapes.geometric, backgrounds, calc, trees, fit}

% Encabezados
\usepackage{fancyhdr}
\pagestyle{fancy}
\fancyhf{}
\fancyhead[L]{\small\itshape UNSAAC -- Ingeniería de Software I}
\fancyhead[R]{\small\itshape SIGC-CUSCO -- V2 Requerimientos y Diseño}
\fancyfoot[C]{\thepage}
\renewcommand{\headrulewidth}{0.4pt}

\usepackage{xurl}
\usepackage{hyperref}
\hypersetup{
    colorlinks=true,
    linkcolor=blue!75!black,
    citecolor=blue!75!black,
    urlcolor=blue!75!black,
    pdftitle={SIGC-CUSCO V2 - Requerimientos y Diseño Arquitectónico},
    pdfauthor={Choquenaira, Yaranga, Ccama - UNSAAC}
}

\usepackage{microtype}
\setlength{\parindent}{0.75cm}
\setlength{\parskip}{0.25em}

\definecolor{SIGCBlue}{HTML}{174A6E}
\definecolor{BoxHeader}{HTML}{E8F1F5}
\definecolor{CodeBg}{HTML}{F6F8FA}

\newcolumntype{Y}{>{\raggedright\arraybackslash}X}
\newcolumntype{P}[1]{>{\raggedright\arraybackslash}p{#1}}

\begin{document}
\hypersetup{pageanchor=false}

% ==============================================================================
% PORTADA INSTITUCIONAL
% ==============================================================================
\begin{titlepage}
    \centering
    \vspace*{0.2cm}
    
    {\fontsize{13}{16}\selectfont \textbf{UNIVERSIDAD NACIONAL DE SAN ANTONIO ABAD DEL CUSCO}}\\[0.2cm]
    {\fontsize{11}{14}\selectfont \textbf{FACULTAD DE INGENIERÍA ELÉCTRICA, ELECTRÓNICA, INFORMÁTICA Y MECÁNICA}}\\[0.2cm]
    {\fontsize{11}{14}\selectfont \textbf{ESCUELA PROFESIONAL DE INGENIERÍA INFORMÁTICA Y DE SISTEMAS}}\\[0.4cm]
    
    \includegraphics[width=3.4cm]{escudo.png}\\[0.5cm]
    
    {\fontsize{13}{16}\selectfont \textbf{DOCUMENTO DE ESPECIFICACIÓN Y DISEÑO ARQUITECTÓNICO (V2)}}\\[0.15cm]
    {\fontsize{10.5}{13.5}\selectfont \textit{Resolución de Ejercicios Ágiles, Requerimientos IEEE 830, Modelado UML y Base de Datos Relacional}}\\[0.25cm]
    \rule{\textwidth}{1.2pt}\\[0.2cm]
    {\fontsize{13}{16.5}\selectfont \textbf{“SISTEMA INTEGRAL DE GESTIÓN DE CAPACITACIONES (SIGC-CUSCO)”}}\\[0.15cm]
    {\fontsize{11}{13}\selectfont \textbf{Incremento Iterativo -- Versión V2.0}}\\[0.15cm]
    \rule{\textwidth}{1.2pt}\\[0.5cm]
    
    \begin{minipage}{0.92\textwidth}
        \setstretch{1.2}
        \textbf{Asignatura:} INGENIERÍA DE SOFTWARE I\\[0.2cm]
        \textbf{Docente:} Ing. Lisha Sabah Diaz Caceres\\[0.2cm]
        \textbf{Enfoque:} Desarrollo Incremental e Iterativo -- Ingeniería de Requisitos\\[0.2cm]
        \textbf{Presentado por:}\\[0.15cm]
        \hspace*{0.4cm}\textbullet\quad Choquenaira Quispe, Noe Franklin \hfill 133962\\[0.12cm]
        \hspace*{0.4cm}\textbullet\quad Yaranga Achahui, Aldo \hfill 103179\\[0.12cm]
        \hspace*{0.4cm}\textbullet\quad Ccama Enriquez, Carolay \hfill 210921
    \end{minipage}\\[0.6cm]
    
    \vfill
    {\textbf{Cusco -- Perú -- 2026}}\\[0.2cm]
\end{titlepage}

% ==============================================================================
% ÍNDICE GENERAL
% ==============================================================================
\newpage
\hypersetup{pageanchor=true}
\pagenumbering{roman}
{
    \hypersetup{linkcolor=black}
    \begin{spacing}{1.0}
    \tableofcontents
    \end{spacing}
}
\newpage
\pagenumbering{arabic}

% ==============================================================================
% SECCIÓN 1: RESOLUCIÓN FORMAL DE EJERCICIOS DE CLASE
% ==============================================================================
\section{Resolución Formal de Ejercicios Prácticos de Flujo Ágil}

A continuación, se presentan las soluciones técnicas y metodológicas a los dos escenarios prácticos planteados en clase, aplicadas rigurosamente al contexto del proyecto \textbf{SIGC-CUSCO}.

\subsection{Ejercicio 1: Una Revisión Saturada (Gestión de Flujo y Cuellos de Botella)}
\textbf{Escenario Planteado:} \textit{El límite de la columna Revisión (Review) es 1 y ya existe un ítem esperando en cola.}

\begin{enumerate}
    \item \textbf{¿Conviene iniciar la tarea de recordatorios? Justificación:}
    \begin{itemize}
        \item \textbf{Decisión Técnica:} \textbf{NO} conviene iniciarla bajo ninguna circunstancia.
        \item \textbf{Fundamento en la Ley de Little y Principios Kanban:} En la teoría de colas y el marco Kanban rige la Ley de Little:
        \begin{equation}
            WIP = TH \times CT \quad \Longrightarrow \quad CT = \frac{WIP}{TH}
        \end{equation}
        donde $WIP$ es el trabajo en curso, $TH$ es el rendimiento (\textit{Throughput}) y $CT$ es el tiempo de ciclo (\textit{Cycle Time}). Cuando la fase de \textit{Review} está saturada ($WIP = 1$), abrir una nueva tarea (como "Módulo de Recordatorios") eleva el $WIP$ global sin incrementar la tasa de entrega ($TH$), disparando el tiempo de ciclo y generando inventario intermedio no verificado. Se aplica el principio rector: \textit{“Stop starting, start finishing”} (Dejar de empezar y enfocarse en terminar).
    \end{itemize}

    \item \textbf{Acuerdo de Colaboración: Estrategia de Swarming / Mob Review:}
    \begin{itemize}
        \item En lugar de abrir código nuevo, el desarrollador disponible aplica \textbf{Swarming} (trabajo concentrado en equipo). Se pone a disposición del revisor para realizar pruebas cruzadas locales (\textit{peer testing}), comprobar los criterios de aceptación y corregir observaciones de código en el acto.
        \item De esta manera, el ítem en revisión se aprueba y se mueve inmediatamente a \textbf{Done}, liberando la capacidad de la columna \textit{Review} para admitir el siguiente elemento del flujo.
    \end{itemize}

    \item \textbf{Protocolo de Actuación ante una Tarjeta Bloqueada:}
    \begin{itemize}
        \item \textbf{Identificación Visual Inmediata:} La tarjeta recibe la etiqueta roja \texttt{[BLOQUEADA]} y se mueve al carril superior del tablero.
        \item \textbf{Registro de Dependencia Externa:} Se anota en la tarjeta la causa raíz específica (ej. \textit{“Pendiente respuesta de API RENIEC para validación de DNI”} o \textit{“Falta plantilla oficial de diploma municipal”}).
        \item \textbf{Política de Aislamiento y SLA:} La tarjeta bloqueada no satura el límite de trabajo activo del desarrollador asignado, pero adquiere prioridad absoluta para el Facilitador/Scrum Master, con un acuerdo de nivel de servicio (SLA) de resolución inferior a 24 horas gestionado en el \textit{Daily Standup}.
    \end{itemize}
\end{enumerate}

\begin{tcolorbox}[colback=blue!4!white, colframe=blue!70!black, title=\textbf{Políticas Explícitas de Flujo en el Tablero SIGC-CUSCO}, arc=1.5mm]
\begin{itemize}\setlength\itemsep{0.15em}
    \item \textbf{Límites WIP Formales:} \textit{To Do} = 5 máx. \textbar{} \textit{In Progress} = 2 máx. \textbar{} \textit{Review / QA} = 1 estricto \textbar{} \textit{Done} = Ilimitado.
    \item \textbf{Regla de Oro ante Bloqueos:} Si una tarjeta se bloquea por más de 12 horas, se destraba de manera colaborativa durante los primeros 5 minutos del \textit{Daily Standup}.
\end{itemize}
\end{tcolorbox}

\subsection{Ejercicio 2: Coordinar una Entrega (Objetivo de Ciclo y Roles)}
\textbf{Escenario Planteado:} \textit{El equipo debe demostrar una funcionalidad utilizable en una entrega incremental.}

\begin{enumerate}
    \item \textbf{Objetivo de Ciclo Orientado a Valor de Negocio (\textit{Sprint Goal}):}
    \begin{quote}
        \textit{“Habilitar que un participante en Cusco complete exitosamente su inscripción digital en línea a un curso municipal o turístico mediante validación de su DNI en tiempo real, obteniendo en pantalla su credencial digital con código QR funcional para el control presencial de asistencia.”}
    \end{quote}
    \textit{Criterio de Valor:} Representa un flujo vertical completo (\textit{thin vertical slice}) directamente operable por el usuario final, no una mera tarea técnica horizontal ("crear tablas o diseñar vistas").

    \item \textbf{Asignación de Roles Operativos en el Equipo:}
    \begin{itemize}\setlength\itemsep{0.2em}
        \item \textbf{Product Owner (Aldo Yaranga):} Responsable del valor de negocio. Define y prioriza los ítems del backlog según las necesidades de municipalidades y agencias turísticas de Cusco; valida el cumplimiento de los criterios de aceptación.
        \item \textbf{Scrum Master / Agile Facilitator (Carolay Ccama):} Responsable del flujo de trabajo y la agilidad del equipo. Vigila el respeto riguroso a los límites WIP en Trello, remueve impedimentos externos y facilita las ceremonias de revisión.
        \item \textbf{Tech Lead / Dev Lead (Noe Franklin Choquenaira):} Responsable de la arquitectura y la calidad del código. Diseña el esquema de datos en PostgreSQL, construye los servicios de backend, integra la seguridad criptográfica y valida las pruebas unitarias.
    \end{itemize}

    \item \textbf{Dos Evidencias Objetivas de Aceptación (\textit{Definition of Done - DoD}):}
    \begin{itemize}\setlength\itemsep{0.2em}
        \item \textbf{Evidencia 1 (Demostración Funcional en Vivo):} Prueba integral en navegador donde se ingresa un DNI, el sistema valida la disponibilidad de vacantes, almacena la inscripción y renderiza de forma instantánea el código QR verificable en pantalla.
        \item \textbf{Evidencia 2 (Garantía de Calidad de Software):} \textit{Pull Request} en GitHub revisado y aprobado por dos integrantes del equipo, con 0 advertencias en análisis estático de código (\textit{linter}) y script SQL ejecutado sin errores de integridad.
    \end{itemize}
\end{enumerate}

% ==============================================================================
% SECCIÓN 2: ESPECIFICACIÓN FORMAL DE REQUERIMIENTOS (IEEE 830 / ISO 29148)
% ==============================================================================
\section{Especificación Formal de Requerimientos de Software}

Bajo los estándares internacionales IEEE 830 e ISO/IEC/IEEE 29148, se formaliza la especificación técnica del sistema \textbf{SIGC-CUSCO}.

\subsection{Requerimientos Funcionales (RF) Priorizados bajo MoSCoW}

\begin{table}[H]
\centering
\small
\renewcommand{\arraystretch}{1.2}
\begin{tabularx}{\textwidth}{P{1.3cm}P{4.0cm}YP{2.2cm}}
\toprule
\textbf{Código} & \textbf{Nombre del Requisito} & \textbf{Especificación Funcional Detallada} & \textbf{Prioridad} \\
\midrule
\textbf{RF-01} & Control de Acceso RBAC & Autenticación segura y autorización por roles: Administrador, Coordinador Municipal/Turístico, Docente y Participante. & Must Have \\
\midrule
\textbf{RF-02} & Gestión Multientidad & Parametrización aislada de instituciones organizadoras (Municipalidades y Empresas de Turismo) con logos oficiales. & Must Have \\
\midrule
\textbf{RF-03} & Catálogo de Capacitaciones & Creación de cursos con definición de fechas, horas pedagógicas, cupo máximo, docente asignado y temario. & Must Have \\
\midrule
\textbf{RF-04} & Inscripción con DNI & Formulario público con validación algorítmica de DNI peruano, verificación de unicidad y control automático de vacantes. & Must Have \\
\midrule
\textbf{RF-05} & Asistencia Digital QR & Registro instantáneo de asistencia mediante escaneo del código QR del participante desde el dispositivo móvil del docente. & Must Have \\
\midrule
\textbf{RF-06} & Calificaciones y Actas & Registro de notas en escala vigesimal (0 a 20), cálculo automático de condición de aprobación y cierre inmutable de actas. & Should Have \\
\midrule
\textbf{RF-07} & Emisión de Certificados & Generación masiva de certificados PDF firmados digitalmente, provistos de código QR y hash de autenticidad SHA-256. & Must Have \\
\midrule
\textbf{RF-08} & Verificador Público Web & Portal web público donde cualquier entidad puede escanear el QR o ingresar el código para validar la autenticidad del diploma. & Must Have \\
\midrule
\textbf{RF-09} & Modo Offline de Asistencia & La interfaz del docente permite registrar asistencias sin conectividad a internet, sincronizando datos al recuperar enlace. & Should Have \\
\midrule
\textbf{RF-10} & Reportes de Fiscalización & Generación de padrones de asistencia y notas exportables a formato Excel (.xlsx) compatibles con auditorías de OSCE. & Could Have \\
\bottomrule
\end{tabularx}
\caption{Matriz de Requerimientos Funcionales con priorización MoSCoW.}
\label{tab:rf}
\end{table}

\subsection{Requerimientos No Funcionales (RNF) según ISO/IEC 25010}

\begin{table}[H]
\centering
\small
\renewcommand{\arraystretch}{1.2}
\begin{tabularx}{\textwidth}{P{1.5cm}P{3.4cm}Y}
\toprule
\textbf{Código} & \textbf{Dimensión de Calidad} & \textbf{Criterio de Aceptación Cuantitativo y Tecnológico} \\
\midrule
\textbf{RNF-01} & \textbf{Seguridad Criptográfica} & Integridad documental mediante hash SHA-256 inmutable por certificado. Almacenamiento de credenciales con cifrado adaptativo \textit{bcrypt} ($cost \ge 12$). \\
\midrule
\textbf{RNF-02} & \textbf{Rendimiento en Lote} & Generación asíncrona de 100 certificados PDF en un tiempo total menor a 15 segundos sin degradar el tiempo de respuesta del servidor web ($< 500$ ms). \\
\midrule
\textbf{RNF-03} & \textbf{Disponibilidad del Servicio} & Arquitectura desplegada en nube con disponibilidad operativa mínima del 99.5\,\% en horarios hábiles de oficina. \\
\midrule
\textbf{RNF-04} & \textbf{Usabilidad Móvil (Responsive)} & Interfaz adaptable que permite completar la inscripción ciudadana en máximo 3 pantallas y menos de 60 segundos en dispositivos móviles. \\
\midrule
\textbf{RNF-05} & \textbf{Tolerancia a Desconexión} & Persistencia local de registros de asistencia mediante \textit{IndexedDB} en el navegador para tolerar cortes de red en locales comunales de Cusco. \\
\bottomrule
\end{tabularx}
\caption{Requerimientos No Funcionales clasificados según el estándar ISO/IEC 25010.}
\label{tab:rnf}
\end{table}

% ==============================================================================
% SECCIÓN 3: MODELADO VISUAL DEL SISTEMA (UML)
% ==============================================================================
\section{Modelado Visual del Sistema (UML)}

\subsection{Diagrama de Casos de Uso del Sistema (UML)}
El modelo de casos de uso (Figura~\ref{fig:casos_uso}) define los límites del sistema \textbf{SIGC-CUSCO}, diferenciando los actores primarios de los servicios de auditoría y especificando las relaciones de inclusión ($\ll\text{include}\gg$).

\begin{figure}[H]
\centering
\begin{tikzpicture}[
    scale=0.92, transform shape,
    actor/.style={rectangle, draw=blue!75!black, fill=blue!10, rounded corners=4pt, font=\footnotesize\bfseries, minimum width=2.6cm, minimum height=0.75cm, align=center},
    usecase/.style={ellipse, draw=black!85, fill=gray!8, font=\scriptsize, align=center, text width=3.0cm, minimum height=0.75cm},
    rel/.style={draw=black!75, -Latex, thick},
    incl/.style={draw=blue!75!black, -Latex, dashed, thick}
]

% Límite del Sistema (System Boundary)
\draw[draw=blue!80!black, thick, fill=blue!1!white, rounded corners=6pt] (-2.2, -4.8) rectangle (6.8, 4.3);
\node[font=\footnotesize\bfseries\color{blue!80!black}, anchor=north west] at (-2.0, 4.1) {Límite del Sistema: SIGC-CUSCO};

% Actores Primarios (Izquierda)
\node[actor] (admin) at (-4.8, 2.8) {Coordinador Edil /\\Turístico};
\node[actor] (docente) at (-4.8, 0.2) {Docente /\\Capacitador};
\node[actor] (alumno) at (-4.8, -2.6) {Participante /\\Ciudadano};

% Actor Secundario (Derecha)
\node[actor] (auditor) at (9.4, -0.8) {Verificador /\\Auditor Público};

% Casos de Uso Principales
\node[usecase] (uc1) at (0.3, 3.2) {Parametrizar Capacitación\\y Asignar Cupos};
\node[usecase] (uc2) at (0.3, 1.8) {Gestionar Inscripciones\\y Vacantes};
\node[usecase] (uc3) at (0.3, 0.4) {Controlar Asistencia\\Digital (QR / Offline)};
\node[usecase] (uc4) at (0.3, -1.0) {Registrar Evaluaciones\\y Cerrar Actas};
\node[usecase] (uc5) at (0.3, -2.4) {Inscribirse en Línea\\a Capacitación};
\node[usecase] (uc6) at (0.3, -3.8) {Descargar Certificado\\Digital Auténtico};

% Casos de Uso Incluidos <<include>>
\node[usecase] (uc7) at (4.7, -2.4) {Validar Formato\\y Unicidad de DNI};
\node[usecase] (uc8) at (4.7, -0.8) {Verificar Hash SHA-256\\y Validez Pública};

% Conexiones de Actores a Casos de Uso
\draw[rel] (admin) -- (uc1);
\draw[rel] (admin) -- (uc2);
\draw[rel] (docente) -- (uc3);
\draw[rel] (docente) -- (uc4);
\draw[rel] (alumno) -- (uc5);
\draw[rel] (alumno) -- (uc6);
\draw[rel] (auditor) -- (uc8);

% Inclusiones
\draw[incl] (uc5) -- node[above, font=\tiny\bfseries] {$\ll$include$\gg$} (uc7);
\draw[incl] (uc6) -- node[above, font=\tiny\bfseries, sloped] {$\ll$include$\gg$} (uc8);

\end{tikzpicture}
\caption{Diagrama de Casos de Uso del Sistema SIGC-CUSCO con delimitación de frontera funcional.}
\label{fig:casos_uso}
\end{figure}

\subsection{Diagrama de Clases del Dominio y Multiplicidades}
El modelo de clases del dominio (Figura~\ref{fig:clases}) representa las entidades estructurales del negocio y sus multiplicidades formales.

\begin{figure}[H]
\centering
\begin{tikzpicture}[
    scale=0.92, transform shape,
    classbox/.style={rectangle, draw=blue!85!black, fill=blue!4!white, rounded corners=2pt, font=\scriptsize, align=left, text width=4.3cm, inner sep=4pt},
    card/.style={font=\tiny\bfseries, text=blue!90!black},
    link/.style={draw=black!80, thick}
]

% Nodos de Clases
\node[classbox] (inst) at (-4.8, 3.4) {
\textbf{\centering Institucion}\\[0.08cm]
\hrule\vspace{0.08cm}
- id\_institucion: int [PK]\\
- ruc: varchar(11) [UQ]\\
- razon\_social: varchar(150)\\
- tipo: enum(Publica, Privada)\\
- logo\_url: text\\[0.08cm]
\hrule\vspace{0.08cm}
+ crear\_curso()\\
+ emitir\_reporte\_auditoria()
};

\node[classbox] (cap) at (1.2, 3.4) {
\textbf{\centering Capacitacion}\\[0.08cm]
\hrule\vspace{0.08cm}
- id\_capacitacion: int [PK]\\
- id\_institucion: int [FK]\\
- id\_docente: int [FK]\\
- titulo: varchar(150)\\
- horas\_academicas: int\\
- vacantes\_max: int\\[0.08cm]
\hrule\vspace{0.08cm}
+ validar\_disponibilidad(): bool\\
+ cerrar\_acta\_final()
};

\node[classbox] (user) at (-4.8, -1.8) {
\textbf{\centering Usuario}\\[0.08cm]
\hrule\vspace{0.08cm}
- id\_usuario: int [PK]\\
- dni: varchar(8) [UQ]\\
- nombres: varchar(80)\\
- correo: varchar(100) [UQ]\\
- rol: enum(Admin, Doc, Part)\\[0.08cm]
\hrule\vspace{0.08cm}
+ autenticar()\\
+ consultar\_diplomas()
};

\node[classbox] (insc) at (1.2, -1.8) {
\textbf{\centering Inscripcion}\\[0.08cm]
\hrule\vspace{0.08cm}
- id\_inscripcion: int [PK]\\
- id\_capacitacion: int [FK]\\
- id\_participante: int [FK]\\
- fecha\_registro: timestamp\\
- estado: varchar(20)\\[0.08cm]
\hrule\vspace{0.08cm}
+ calcular\_asistencias(): float\\
+ verificar\_aprobacion(): bool
};

\node[classbox] (cert) at (7.0, -1.8) {
\textbf{\centering Certificado}\\[0.08cm]
\hrule\vspace{0.08cm}
- id\_certificado: int [PK]\\
- id\_inscripcion: int [FK, UQ]\\
- codigo\_unico: varchar(32) [UQ]\\
- hash\_sha256: char(64)\\
- url\_pdf: text\\[0.08cm]
\hrule\vspace{0.08cm}
+ generar\_pdf\_qr()\\
+ verificar\_autenticidad(): bool
};

% Conexiones con multiplicidades
\draw[link] (inst) -- node[above, card] {1} node[above, card, pos=0.88] {1..*} (cap);
\draw[link] (user) -- node[above, card] {1} node[above, card, pos=0.88] {0..*} (insc);
\draw[link] (cap) -- node[right, card] {1} node[right, card, pos=0.88] {0..*} (insc);
\draw[link] (insc) -- node[above, card] {1} node[above, card, pos=0.88] {0..1} (cert);

\end{tikzpicture}
\caption{Diagrama de Clases del Dominio con atributos, métodos y cardinalidades normalizadas.}
\label{fig:clases}
\end{figure}

% ==============================================================================
% SECCIÓN 4: BASE DE DATOS Y ESQUEMA RELACIONAL (POSTGRESQL)
% ==============================================================================
\section{Diseño de Base de Datos Relacional y Esquema Físico}

El sistema implementa el gestor relacional \textbf{PostgreSQL 16} para asegurar soporte estricto de transacciones ACID y alta consistencia en la emisión documental.

\subsection{Modelo Relacional y Reglas de Integridad Referencial}
El diseño de datos responde a la Tercera Forma Normal (3FN), garantizando que no existan redundancias ni anomalías de actualización:
\begin{enumerate}\setlength\itemsep{0.2em}
    \item \textbf{Relación Uno a Muchos (1 : N) entre Institución y Capacitaciones:} Una entidad (municipalidad o empresa) patrocina múltiples cursos a lo largo del año. La eliminación de una institución se restringe si tiene cursos históricos cerrados (\texttt{ON DELETE RESTRICT}) para preservar la trazabilidad de auditoría.
    \item \textbf{Relación Muchos a Muchos (N : M) resuelta por Entidad Asociativa:} Un participante se matricula en múltiples capacitaciones y una capacitación alberga decenas de alumnos. Esta asociación se materializa en la tabla \textbf{\texttt{inscripciones}}, la cual incluye una restricción única compuesta \texttt{UNIQUE (id\_capacitacion, id\_participante)} para impedir duplicidades de matrícula.
    \item \textbf{Relación Uno a Uno Condicional (1 : 0..1) con Certificados:} Una inscripción completada y aprobada genera a lo sumo un único certificado oficial. Se implementa con una restricción \texttt{UNIQUE} sobre la clave foránea \texttt{id\_inscripcion} en la tabla \texttt{certificados}.
\end{enumerate}

\subsection{Script DDL de Implementación Física (PostgreSQL)}

\begin{tcolorbox}[breakable, colback=CodeBg, colframe=black!40, title=\textbf{Script DDL PostgreSQL 16 -- SIGC-CUSCO}, arc=1mm]
\scriptsize
\begin{verbatim}
-- 1. Entidades Organizadoras (Multitenancy)
CREATE TABLE instituciones (
    id_institucion SERIAL PRIMARY KEY,
    ruc VARCHAR(11) UNIQUE NOT NULL,
    razon_social VARCHAR(150) NOT NULL,
    tipo VARCHAR(20) NOT NULL CHECK (tipo IN ('Publica', 'Privada')),
    logo_url TEXT,
    activo BOOLEAN DEFAULT TRUE,
    creado_en TIMESTAMP DEFAULT CURRENT_TIMESTAMP
);

-- 2. Usuarios del Sistema y Roles RBAC
CREATE TABLE usuarios (
    id_usuario SERIAL PRIMARY KEY,
    dni VARCHAR(8) UNIQUE NOT NULL,
    nombres VARCHAR(80) NOT NULL,
    apellidos VARCHAR(80) NOT NULL,
    correo VARCHAR(100) UNIQUE NOT NULL,
    password_hash VARCHAR(255) NOT NULL,
    rol VARCHAR(20) NOT NULL CHECK (rol IN ('Admin', 'Coordinador', 'Docente', 'Alumno')),
    creado_en TIMESTAMP DEFAULT CURRENT_TIMESTAMP
);

-- 3. Catálogo de Capacitaciones
CREATE TABLE capacitaciones (
    id_capacitacion SERIAL PRIMARY KEY,
    id_institucion INT NOT NULL 
        REFERENCES instituciones(id_institucion) ON DELETE RESTRICT,
    id_docente INT 
        REFERENCES usuarios(id_usuario) ON DELETE SET NULL,
    titulo VARCHAR(150) NOT NULL,
    horas_academicas INT NOT NULL CHECK (horas_academicas > 0),
    vacantes_max INT NOT NULL CHECK (vacantes_max > 0),
    fecha_inicio DATE NOT NULL,
    fecha_fin DATE,
    estado VARCHAR(20) DEFAULT 'Planificada' 
        CHECK (estado IN ('Planificada', 'En_Curso', 'Finalizada'))
);

-- 4. Matrículas e Inscripciones (Tabla Intermedia)
CREATE TABLE inscripciones (
    id_inscripcion SERIAL PRIMARY KEY,
    id_capacitacion INT NOT NULL 
        REFERENCES capacitaciones(id_capacitacion) ON DELETE CASCADE,
    id_participante INT NOT NULL 
        REFERENCES usuarios(id_usuario) ON DELETE RESTRICT,
    fecha_registro TIMESTAMP DEFAULT CURRENT_TIMESTAMP,
    nota_final NUMERIC(4,2) CHECK (nota_final >= 0 AND nota_final <= 20),
    estado VARCHAR(20) DEFAULT 'Confirmada' 
        CHECK (estado IN ('Confirmada', 'Aprobado', 'Desaprobado', 'Anulada')),
    CONSTRAINT uq_participante_curso UNIQUE (id_capacitacion, id_participante)
);

-- 5. Control de Asistencias por Sesión (Soporte QR / Offline)
CREATE TABLE asistencias (
    id_asistencia SERIAL PRIMARY KEY,
    id_inscripcion INT NOT NULL 
        REFERENCES inscripciones(id_inscripcion) ON DELETE CASCADE,
    numero_sesion INT NOT NULL CHECK (numero_sesion > 0),
    fecha_hora TIMESTAMP DEFAULT CURRENT_TIMESTAMP,
    estado VARCHAR(15) NOT NULL CHECK (estado IN ('Presente', 'Tardanza', 'Falta')),
    metodo VARCHAR(15) DEFAULT 'QR_Digital' 
        CHECK (metodo IN ('QR_Digital', 'PIN_Manual', 'Offline_Sync'))
);

-- 6. Emisión Criptográfica de Certificados
CREATE TABLE certificados (
    id_certificado SERIAL PRIMARY KEY,
    id_inscripcion INT UNIQUE NOT NULL 
        REFERENCES inscripciones(id_inscripcion) ON DELETE RESTRICT,
    codigo_unico VARCHAR(32) UNIQUE NOT NULL,
    hash_sha256 CHAR(64) NOT NULL,
    fecha_emision TIMESTAMP DEFAULT CURRENT_TIMESTAMP,
    url_pdf TEXT NOT NULL
);

-- Índices de Alto Rendimiento para Búsqueda y Validación
CREATE INDEX idx_usuarios_dni ON usuarios(dni);
CREATE INDEX idx_inscripciones_curso ON inscripciones(id_capacitacion);
CREATE INDEX idx_certificados_hash ON certificados(hash_sha256);
CREATE INDEX idx_certificados_codigo ON certificados(codigo_unico);
\end{verbatim}
\end{tcolorbox}

% ==============================================================================
% SECCIÓN 5: ARQUITECTURA TECNOLÓGICA Y STACK DEL PROYECTO
% ==============================================================================
\section{Arquitectura de Software y Stack Tecnológico del Proyecto}

Para garantizar robustez, mantenibilidad y escalabilidad, el proyecto \textbf{SIGC-CUSCO} adopta una arquitectura en tres capas desacopladas con componentes orientados a servicios RESTful.

\subsection{Capas de la Arquitectura de Software}
\begin{enumerate}\setlength\itemsep{0.2em}
    \item \textbf{Capa de Presentación (Frontend Web):} Desarrollada en \textbf{React 18} con empaquetador \textbf{Vite} y estilos utilitarios en \textbf{Tailwind CSS}. Ofrece una interfaz ligera, accesible desde teléfonos móviles de gama baja y optimizada para la interacción rápida de ciudadanos y docentes.
    \item \textbf{Capa de Lógica de Negocio (Backend API):} Implementada con \textbf{Node.js} y \textbf{Express} bajo el estándar tipado de \textbf{TypeScript}. Gestiona la validación algorítmica de DNIs, el control concurrente de cupos, la orquestación del marcado de asistencia y la cola asíncrona de procesamiento de PDFs.
    \item \textbf{Capa de Persistencia y Seguridad (Datos):} Motor \textbf{PostgreSQL 16} para el almacenamiento relacional estructurado. Generación de códigos QR mediante librerías nativas vectoriales y cálculo de resumen criptográfico inmutable mediante el algoritmo estándar SHA-256.
\end{enumerate}

\subsection{Resumen de Tecnologías y Entorno de Desarrollo}

\begin{table}[H]
\centering
\small
\renewcommand{\arraystretch}{1.2}
\begin{tabularx}{\textwidth}{P{3.0cm}P{4.2cm}Y}
\toprule
\textbf{Componente / Capa} & \textbf{Tecnología Elegida} & \textbf{Justificación Técnica en el Proyecto} \\
\midrule
\textbf{Base de Datos} & PostgreSQL 16 & Soporte ACID estricto, integridad referencial y alto rendimiento en consultas indexadas. \\
\midrule
\textbf{Backend / API} & Node.js + Express (TypeScript) & Asincronía eficiente (\textit{Event Loop}) para procesar inscripciones masivas sin bloqueos. \\
\midrule
\textbf{Frontend Web} & React + Vite + Tailwind CSS & Renderizado rápido, diseño 100\,\% adaptable para smartphones y soporte de PWA. \\
\midrule
\textbf{Seguridad y Hash} & Algoritmo SHA-256 / Bcrypt & Certificados infalsificables auditables y protección robusta de contraseñas. \\
\midrule
\textbf{Control de Versiones} & Git \& GitHub (\texttt{SIGC-CUSCO}) & Colaboración concurrente, trazabilidad de código y despliegue continuo (CI/CD). \\
\midrule
\textbf{Gestión Ágil} & Tablero Kanban en Trello & Visualización de flujo, políticas de bloqueo y control de cuellos de botella con límites WIP. \\
\bottomrule
\end{tabularx}
\caption{Stack tecnológico formal adoptado para la construcción de SIGC-CUSCO.}
\label{tab:stack}
\end{table}

% ==============================================================================
% SECCIÓN 6: PLAN DE EVOLUCIÓN INCREMENTAL E ITERATIVO
% ==============================================================================
\section{Plan de Evolución Incremental e Iterativo por Sprints}

Para cumplir con la entrega continua de valor en cada sesión del curso, el desarrollo de \textbf{SIGC-CUSCO} se planifica en cinco incrementos progresivos:

\begin{table}[H]
\centering
\small
\renewcommand{\arraystretch}{1.2}
\begin{tabularx}{\textwidth}{P{2.6cm}P{4.4cm}Y}
\toprule
\textbf{Incremento / Sprint} & \textbf{Entregable Técnico Principal} & \textbf{Criterio de Validación Objetiva} \\
\midrule
\textbf{Sprint 1: Inception} & Visión, Lean Canvas y Tablero Trello. & Aprobación de alcance y matriz de stakeholders en repositorio GitHub. \\
\midrule
\textbf{Sprint 2: Requisitos y Arquitectura (Actual)} & Casos de uso, clases de dominio, DER y script DDL en PostgreSQL. & Compilación formal del documento V2 y ejecución limpia del esquema SQL. \\
\midrule
\textbf{Sprint 3: Core de Inscripciones} & Portal público con validación de DNI y gestión de cupos en tiempo real. & Demostración de registro ciudadano y persistencia en base de datos. \\
\midrule
\textbf{Sprint 4: Asistencia Digital} & Lector móvil de código QR con persistencia offline en IndexedDB. & Marcado de presencia en aula comunal y sincronización con el servidor. \\
\midrule
\textbf{Sprint 5: Certificación y Auditoría} & Motor de emisión de diplomas PDF con hash SHA-256 y validador público web. & Verificación exitosa de autenticidad escaneando el código QR con cualquier celular. \\
\bottomrule
\end{tabularx}
\caption{Roadmap de desarrollo incremental e iterativo del sistema SIGC-CUSCO.}
\label{tab:roadmap}
\end{table}

\end{document}
